%========================================
% Industrial Training Logbook (LaTeX)
% Student: Nada Ashraf Moussa Kamel Gomaa
% ID: 120210358
% Company: CLS - Learning Solutions (under DEPI)
% Training Period: 7 Months (April–October 2024)
%========================================

\documentclass[12pt,a4paper]{article}

\usepackage[margin=1.5cm]{geometry}
\usepackage{array}
\usepackage{tabularx}
\usepackage{enumitem}
\usepackage{hyperref}
\usepackage{xcolor}
\usepackage{tcolorbox}
\usepackage{setspace}
\usepackage{graphicx}

\hypersetup{
  colorlinks=true,
  linkcolor=blue,
  urlcolor=blue
}

\setlength{\parindent}{0pt}
\setstretch{1.1}

% ---------- Styles ----------
\definecolor{boxgray}{RGB}{245,245,245}
\definecolor{linegray}{RGB}{190,190,190}

\newcolumntype{Y}{>{\raggedright\arraybackslash}X}

% ---------- Week Entry Macro ----------
% #1 Date Range
% #2 Week No(s)
% #3 Activities
% #4 Summary
% #5 Key Topics
% #6 Achievement of the Day
% #7 Functional Skills (bullets)
% #8 Soft Skills (bullets)
\newcommand{\WeekEntry}[8]{
\begin{tcolorbox}[
  colback=boxgray,
  colframe=linegray,
  boxrule=0.6pt,
  arc=1.5mm,
  left=1mm,right=1mm,top=2mm,bottom=2mm
]
\noindent
\begin{tabularx}{\linewidth}{|p{2.8cm}|Y|}
\hline
\multicolumn{2}{|l|}{\textbf{Student's Logbook}}\\ \hline
\textbf{Date:} & #1 \\ \hline
\multicolumn{2}{|c|}{\textbf{#2}}\\ \hline
\textbf{Activities:} &
\begin{itemize}[leftmargin=1em,nosep,topsep=0pt,itemsep=1pt]
#3
\end{itemize}
\\ \hline
\textbf{Summary:} & #4 \\ \hline
\textbf{Key Topics:} &
\begin{itemize}[leftmargin=1em,nosep,topsep=0pt,itemsep=1pt]
#5
\end{itemize}
\\ \hline
\multicolumn{2}{|p{\dimexpr\linewidth-2\tabcolsep-2\arrayrulewidth}|}{\textbf{Achievement of the Day:} #6}\\ \hline
\textbf{Functional Skills:} &
\begin{itemize}[leftmargin=1em,nosep,topsep=0pt,itemsep=2pt]
#7
\end{itemize}
\\ \hline
\textbf{Soft Skills:} &
\begin{itemize}[leftmargin=1em,nosep,topsep=0pt,itemsep=2pt]
#8
\end{itemize}
\\ \hline
\end{tabularx}
\end{tcolorbox}
}

\begin{document}

% ---------- Header ----------
{\LARGE \textbf{Industrial Training Logbook Year: 2024}}\\[6pt]

\textbf{Student's Name:} Nada Ashraf Moussa Kamel Gomaa\\
\textbf{Student's ID:} 120210358\\[10pt]

% ---------- Declaration ----------
{\Large \textbf{Declaration}}\\[6pt]
Before submitting this document for assessment, please complete the following declaration:\\[6pt]

This is to certify that the contents of the Industrial Training Logbook are a true and accurate reflection of the work done by the student in the training company.\\[10pt]

\textbf{Student's name:} Nada Ashraf Moussa Kamel Gomaa\\
\textbf{Student's signature:} Nada Ashraf Moussa Kamel Gomaa\\
\textbf{Supervisor's name:} Egypt-Digital-Pioneers-Initiative\\
\textbf{Supervisor's signature:} DEPI- Dr. Eman Raslan\\
\textbf{Company's name:} CLS- Learning Solutions\\[14pt]

% ---------- Student Details ----------
{\Large \textbf{Student's Details}}\\[6pt]
\textbf{Name:} Nada Ashraf Moussa Kamel Gomaa\\
\textbf{I/C No.:} \hspace{2cm} \textbf{Matric No.:}\\
\textbf{Mobile Tel. No.:} 01093287761\\
\textbf{House Tel. No.:} 01093287761\\
\textbf{Email:} \href{mailto:Nada.gomaa@ejust.edu.eg}{Nada.gomaa@ejust.edu.eg}\\
\textbf{Major:} Computer Science and Engineering\\
\textbf{Minor:} ECCE\\[14pt]

% ---------- Placement Details ----------
{\Large \textbf{Placement Details}}\\[6pt]
\textbf{Name of Company:} Egypt Digital Pioneers Initiative (CLS)\\
\textbf{Company Address:} 26 Sharm El-Dokki Main Street, Next to Misr Gas Station in front of Dokki Metro Station, Dokki, Giza Governorate\\
\textbf{Industrial Supervisor:} CLS learning solutions\\
\textbf{Company Contact Tel No.:} +20 11 53615270 \hspace{1.2cm} \textbf{Fax No.:} +20 11 53615270\\
\textbf{Training Period:} 7 Months (April to October 2024)\\
\textbf{Allowance:} --\\
(Indicate the amount of allowance paid per month)\\[18pt]

{\Large \textbf{Student's Logbook}}\\[6pt]
\textbf{Training intensity assumption used for planning:} 3--4 meetings per week (3--4 hours each) + continuous labs and self-study.

%========================================
% WEEKS 1–2
%========================================
\WeekEntry
{1 April -- 14 April 2024}
{Weeks 1--2}
{
\item Orientation to the training structure and deliverables (logbook discipline, weekly reporting, and lab submission workflow).
\item Orientation to Microsoft SQL and relational database fundamentals.
\item Learning SQL essentials: DDL (CREATE/ALTER/DROP) and DML (SELECT/INSERT/UPDATE/DELETE).
\item Hands-on: database setup and environment configuration for practice.
\item Hands-on: writing simple queries with filters, sorting, grouping, and joins.
\item Short exercises in query debugging, fixing syntax/logic errors, and validating outputs using sample data.
}
{
Initial phase focused on SQL foundations and establishing a correct working environment, followed by intensive practice on essential query patterns required for data engineering and analytics workflows.
}
{
\item Relational databases and SQL Server environment setup
\item DDL vs DML and correct usage
\item Core SELECT queries, WHERE, ORDER BY, GROUP BY
\item JOIN types (INNER/LEFT/RIGHT) and practical use cases
\item Query troubleshooting and result validation
}
{
Developed foundational skills in computer usage, problem-solving, teamwork, and communication
}
{
\item Computer Skills
\item Problem-solving Skills
}
{
\item Communication Skills
\item Teamwork Skills
}

%========================================
% WEEKS 3–4
%========================================
\WeekEntry
{15 April -- 28 April 2024}
{Weeks 3--4}
{
\item Advanced SQL practice: multi-table joins, subqueries, correlated queries, and common table expressions (CTEs).
\item Data quality logic in SQL: NULL handling, constraints awareness, and consistency checks.
\item Performance basics: understanding indexes at a conceptual level and why queries slow down.
\item Applied tasks: create tables + insert sample data + write analytical queries that simulate business reporting.
\item Documentation habit: writing query comments, assumptions, and expected outputs for reproducibility.
}
{
Deepened SQL capability from basic querying to more complex analytical and engineering-oriented patterns, emphasizing correctness, validation, and early performance awareness.
}
{
\item CTEs, subqueries, correlated queries
\item Aggregation with HAVING and multi-step analytics
\item Data validation logic (NULL, missing values, duplicates conceptually)
\item Basic query performance considerations and indexing concepts
\item Reproducible SQL documentation practices
}
{
Advanced SQL proficiency including complex queries, CTEs, and performance optimization techniques
}
{
\item Organizing/Analyzing Data
\item Decision-making Ability
}
{
\item Critical Thinking Skills
\item Teamwork Skills
}

%========================================
% WEEKS 5–6
%========================================
\WeekEntry
{29 April -- 12 May 2024}
{Weeks 5--6}
{
\item Implementing SQL Data Warehouse concepts: OLTP vs OLAP and why analytics requires different schema design.
\item Dimensional modeling practice: fact tables, dimension tables, and star schema vs snowflake schema.
\item Normalization vs denormalization trade-offs and when each is justified.
\item Hands-on: sketching a small warehouse schema for analytics reporting scenarios.
\item Designed example KPIs and mapped them to facts/measures and dimensions.
}
{
Shifted from operational SQL usage to analytical system design: learned how warehouse schemas support reporting performance, consistent metrics, and scalable analytics.
}
{
\item Data warehouse fundamentals (OLTP vs OLAP)
\item Fact vs Dimension tables and grain definition
\item Star vs snowflake schema
\item Normalization vs denormalization trade-offs
\item KPI mapping to schema (measures/dimensions)
}
{
Understanding of data warehouse architecture and dimensional modeling for analytics
}
{
\item Organizing/Analyzing Data
\item Decision-making Ability
}
{
\item Teamwork Skills
\item Critical Thinking Skills
}

%========================================
% WEEKS 7–8
%========================================
\WeekEntry
{13 May -- 26 May 2024}
{Weeks 7--8}
{
\item ETL (Extract, Transform, Load) process concepts and why ETL is essential for warehouse reliability.
\item Data pipeline steps: extracting from sources, cleaning, transforming, and loading into target schema.
\item Data integrity and quality checks: consistency rules, duplicate handling, missing data strategies.
\item Practical mini-pipeline: simulate ETL with SQL scripts and structured transformation steps.
\item Team-based review sessions: comparing schema designs, discussing transformation rules and validation logic.
}
{
Built a structured understanding of ETL pipelines and quality controls required to maintain trustworthy analytics datasets in a warehouse environment.
}
{
\item ETL workflow and stages
\item Data quality rules and validation checks
\item Transformations aligned to business definitions
\item Loading strategies and schema alignment
\item Documentation of assumptions and pipeline steps
}
{
Practical ETL pipeline development and data quality management skills
}
{
\item Organizing/Analyzing Data
\item Problem-solving Skills
}
{
\item Communication Skills
\item Teamwork Skills
}

%========================================
% WEEKS 9–10
%========================================
\WeekEntry
{27 May -- 9 June 2024}
{Weeks 9--10}
{
\item Python programming fundamentals for data workflows: syntax, variables, loops, functions.
\item Data structures: lists, tuples, dictionaries, sets; selecting structures based on task needs.
\item Error handling with try/except; writing safer scripts.
\item File handling and directory operations: reading/writing files, working with paths.
\item Initial work with NumPy arrays and basic Pandas operations for data manipulation.
}
{
Established Python fundamentals required for automation, data preprocessing, and future ML experimentation. Emphasis was placed on correctness, readable code, and controlled error handling.
}
{
\item Python basics (functions, loops, conditions)
\item Data structures and practical usage
\item Error handling and debugging workflow
\item File handling and scripting habits
\item NumPy and Pandas basics
}
{
Foundation in Python programming for data workflows and automation
}
{
\item Computer Skills
\item Quantitative \& Analytical Skills
}
{
\item Communication Skills
\item Interpersonal Skills
}

%========================================
% WEEKS 11–12
%========================================
\WeekEntry
{10 June -- 23 June 2024}
{Weeks 11--12}
{
\item Data preprocessing with Python: handling missing values, duplicates conceptually, encoding categorical data (concept-level).
\item Exploratory Data Analysis basics: summary stats, simple plots, and interpretation.
\item Intro to ML workflow idea: data $\rightarrow$ split $\rightarrow$ train $\rightarrow$ evaluate (concept-level).
\item Strengthened notebook discipline: clean sections, comments, and reproducible runs.
\item Group sessions: presenting solutions, peer review of code clarity and correctness.
}
{
Expanded Python skills toward data preparation and analysis patterns that directly support ML and data engineering deliverables.
}
{
\item Data preprocessing patterns in Python
\item Basic EDA workflow and interpretation
\item Reproducible notebook structure and documentation
\item Clean coding practices and peer review
\item Intro ML pipeline framing (conceptual)
}
{
Data preprocessing and exploratory analysis skills for machine learning readiness
}
{
\item Organizing/Analyzing Data
\item Problem-solving Skills
}
{
\item Critical Thinking Skills
\item Teamwork Skills
}

%========================================
% WEEKS 13–14
%========================================
\WeekEntry
{24 June -- 7 July 2024}
{Weeks 13--14}
{
\item Introduction to Agile development and Scrum roles/events/artifacts.
\item Participation in sprint planning and daily stand-up simulations.
\item Understanding tickets/issue tracking systems: breaking tasks, estimating, defining acceptance criteria.
\item Documentation of sprint tasks and aligning technical deliverables with sprint goals.
\item Collaboration practice: team communication rules, meeting notes, and action items.
}
{
Learned Agile/Scrum as an operational framework for managing technical tasks, coordinating teams, and ensuring traceable progress across complex training modules and project milestones.
}
{
\item Agile principles and Scrum framework
\item Sprint planning, daily standups, task breakdown
\item Ticketing workflow and acceptance criteria
\item Documentation habits for engineering projects
\item Collaboration and coordination routines
}
{
Practical ETL pipeline development and data quality management skills
}
{
\item Organizing/Analyzing Data
\item Problem-solving Skills
}
{
\item Communication Skills
\item Teamwork Skills
}

%========================================
% WEEKS 15–16
%========================================
\WeekEntry
{8 July -- 21 July 2024}
{Weeks 15--16}
{
\item Azure Data Fundamentals: core cloud concepts and why cloud matters for data/ML systems.
\item Understanding cloud storage, compute, and security concepts; shared responsibility concept-level.
\item Hands-on with Azure Data Studio: connecting, exploring data, and running queries.
\item Relational vs non-relational databases: use-cases and tradeoffs.
\item Security and compliance awareness for data systems (concept-level).
}
{
Developed baseline cloud literacy focused on data platforms, including tool usage (Azure Data Studio) and fundamental design/security considerations.
}
{
\item Cloud data concepts and service categories
\item Azure Data Studio workflow
\item Relational vs non-relational databases
\item Security/compliance awareness (fundamentals)
\item Cloud value for business and scalability
}
{
Cloud data platform fundamentals and Azure Data Studio proficiency
}
{
\item Computer Skills
\item Decision-making Ability
}
{
\item Communication Skills
\item Interpersonal Skills
}

%========================================
% WEEKS 17–18
%========================================
\WeekEntry
{22 July -- 4 August 2024}
{Weeks 17--18}
{
\item Advanced Azure discussion: designing secure and scalable data solutions (concept-level).
\item Data ingestion and pipeline thinking: batch vs streaming (concept-level).
\item Monitoring and optimization mindset for cloud workflows (concept-level).
\item Team-based case studies: mapping business requirements to data services and pipeline components.
\item Consolidation quiz sessions and practical review.
}
{
Moved from fundamental cloud concepts to solution-thinking: how to select services and plan pipelines while considering scalability, monitoring, and security.
}
{
\item Designing data solutions in Azure (concepts)
\item Batch vs streaming ingestion (overview)
\item Monitoring and optimization mindset
\item Security and scalability considerations
\item Translating requirements to architecture
}
{
Advanced cloud solution design thinking and architecture planning
}
{
\item Organizing/Analyzing Data
\item Decision-making Ability
}
{
\item Leadership Skills
\item Communication Skills
}

%========================================
% WEEKS 19–20 (Course 7: NLP Foundations)
%========================================
\WeekEntry
{5 August -- 18 August 2024}
{Weeks 19--20}
{
\item Course 7: NLP scope and pipeline: text acquisition, cleaning, normalization, tokenization, modeling, evaluation.
\item NLU vs NLG: understanding vs generation tasks and examples.
\item NLP applications covered: toxicity classification, topic classification, grammatical error correction, information retrieval, question answering.
\item Key NLP challenges: ambiguity (lexical/syntactic/semantic), context dependence, noisy web data, multilinguality/low-resource languages, implicit meaning (sarcasm/cultural context).
\item Preprocessing toolkit concepts: lowercase tradeoffs, regex concept practice, tokenization types (character/word/sentence/subword).
\item Tool comparison: NLTK vs spaCy (use-cases, tokenization support, pipeline differences).
}
{
Built the foundations of NLP workflow and challenges, linking preprocessing decisions to downstream model performance and evaluation reliability.
}
{
\item NLP pipeline and task taxonomy (NLU/NLG)
\item Text cleaning and normalization decisions
\item Tokenization types and trade-offs
\item NLTK vs spaCy overview
\item Real-world NLP application mapping
}
{
NLP pipeline understanding and text preprocessing proficiency
}
{
\item Problem-solving Skills
\item Quantitative \& Analytical Skills
}
{
\item Critical Thinking Skills
\item Teamwork Skills
}

%========================================
% WEEKS 21–22 (Course 7: Traditional NLP)
%========================================
\WeekEntry
{19 August -- 1 September 2024}
{Weeks 21--22}
{
\item Feature engineering for text: one-hot encoding (limitations), Bag of Words, TF-IDF, and N-grams.
\item Understanding sparsity/dimensionality problems and why vocabulary control matters.
\item Traditional ML modeling for NLP: logistic regression, naive Bayes, decision trees (concept + workflow).
\item Practice workflow: vectorizer $\rightarrow$ model $\rightarrow$ evaluation (concept-level).
\item Lab-style discussion based on supervised classification pipelines (e.g., fake news-style classification as a practice theme).
}
{
Advanced from preprocessing to classical NLP representations and traditional ML models, emphasizing interpretability, pipeline structure, and evaluation discipline.
}
{
\item One-hot, BoW, TF-IDF, N-grams
\item Sparsity and dimensionality concerns
\item Classical supervised NLP models
\item Vectorizer-to-classifier pipelines
\item Evaluation framing and error analysis mindset
}
{
Classical NLP techniques and feature engineering for text data
}
{
\item Organizing/Analyzing Data
\item Decision-making Ability
}
{
\item Teamwork Skills
\item Critical Thinking Skills
}

%========================================
% WEEKS 23–24 (Course 7: Deep Learning NLP)
%========================================
\WeekEntry
{2 September -- 15 September 2024}
{Weeks 23--24}
{
\item Deep learning NLP evolution: embeddings and sequence modeling.
\item Embeddings: Word2Vec (CBOW/Skip-gram), GloVe, FastText; contextual embeddings concept: ELMo/BERT/GPT/T5.
\item RNN fundamentals and gradient issues; motivation for LSTM and GRU.
\item Seq2Seq concept: encoder/decoder, bottleneck problem, and attention mechanism concept.
\item Evaluation metrics overview: BLEU (translation), ROUGE (summarization), WER (speech).
\item Transformers overview and why attention is important for long-range dependencies.
}
{
Covered deep-learning NLP foundations and the motivation for transformer architectures, linking model families to tasks and evaluation metrics.
}
{
\item Static vs contextual embeddings
\item RNN/LSTM/GRU concept and limitations
\item Seq2Seq and attention concept
\item BLEU/ROUGE/WER evaluation overview
\item Transformer motivation and attention role
}
{
Deep learning NLP foundations and transformer architecture understanding
}
{
\item Quantitative \& Analytical Skills
\item Problem-solving Skills
}
{
\item Leadership Skills
\item Communication Skills
}

%========================================
% WEEKS 25–26 (Course 8: Hugging Face)
%========================================
\WeekEntry
{16 September -- 29 September 2024}
{Weeks 25--26}
{
\item Course 8: Hugging Face overview and ecosystem components: Hub, Transformers, Datasets, Tokenizers, Spaces, Accelerate.
\item Hub usage concepts: model hub, dataset hub, model/dataset cards, hosted widgets for quick testing.
\item Transformers library: pre-trained model families and tasks (classification, generation, QA, translation, NER, summarization).
\item Pipelines abstraction: tokenization + inference + post-processing as an end-to-end workflow.
\item Model architecture families mapping:
  \begin{itemize}[leftmargin=*,nosep]
    \item Encoder: BERT/DistilBERT/RoBERTa for classification/NER/extractive QA.
    \item Decoder: GPT-style for text generation.
    \item Encoder-decoder: T5/BART/mBART/Marian for translation/summarization/generative QA.
  \end{itemize}
\item Discussed transformer challenges: language dominance (English), data availability, long-document cost, opacity, bias in internet-trained models.
}
{
Established working knowledge of Hugging Face tooling and how its libraries and Hub accelerate experimentation, sharing, and deployment for transformer-based workflows.
}
{
\item Hugging Face Hub (models/datasets/spaces) and cards
\item Transformers library and pipelines
\item Datasets and Tokenizers libraries
\item Model family selection by task
\item Known transformer limitations: bias/opacity/long context
}
{
Hugging Face ecosystem proficiency and transformer model selection skills
}
{
\item Computer Skills
\item Organizing/Analyzing Data
}
{
\item Communication Skills
\item Teamwork Skills
}

%========================================
% WEEKS 27–28 (Course 8: MLOps + MLflow)
%========================================
\WeekEntry
{30 September -- 13 October 2024}
{Weeks 27--28}
{
\item MLOps definition: practices to manage ML development, deployment, maintenance with automation across the lifecycle.
\item Why MLOps: scalability, reproducibility, fast time-to-market, and monitoring.
\item Core data governance concepts: data lineage, data provenance, metadata.
\item CI/CD in MLOps: automated testing of code and models, version control for code/data/models, automated data validation.
\item Continuous Training (CT): retraining when data profile changes and drift affects performance.
\item MLflow overview: Tracking, Projects, Models, Model Registry.
\item MLflow Tracking: logging parameters, metrics, artifacts and using UI for comparison.
\item MLflow Projects: packaging code + environment for reproducibility (Conda/Docker concept).
\item MLflow Models: packaging and deployment options (local, REST, cloud/container concept).
\item MLflow Model Registry: versioning + staging + production + archived lifecycle control.
}
{
Built MLOps literacy centered on MLflow as a lifecycle management tool, connecting governance and CI/CD/CT concepts to reliable ML delivery.
}
{
\item MLOps lifecycle and automation rationale
\item Data lineage/provenance/metadata
\item CI/CD adaptations for ML + Continuous Training
\item MLflow Tracking/Projects/Models/Registry
\item Reproducibility and versioning practices
}
{
MLOps practices and MLflow lifecycle management understanding
}
{
\item Quantitative \& Analytical Skills
\item Problem-solving Skills
}
{
\item Critical Thinking Skills
\item Teamwork Skills
}

%========================================
% WEEKS 29–30 (Course 9: Prompt Engineering + FINAL PROJECT)
%========================================
\WeekEntry
{14 October -- 27 October 2024}
{Weeks 29--30}
{
\item Course 9 (Prompt Engineering): prompt structure, constraints, format control, examples (zero-shot vs few-shot), iterative refinement and evaluation prompts.
\item Applied prompt engineering for practical work: debugging assistance, refactoring notebooks, generating checklists and edge cases, improving documentation wording and structure.
\item Final applied project: \textbf{cheXray\_Classifier} (Pneumonia vs Normal) using Kaggle dataset:
  \begin{itemize}[leftmargin=*,nosep]
    \item Dataset: 5,863 pediatric chest X-rays; folders train/test/val with NORMAL/PNEUMONIA.
    \item Context and reference paper link included in resources.
  \end{itemize}
\item Colab pipeline implementation (end-to-end):
  \begin{itemize}[leftmargin=*,nosep]
    \item Kaggle API setup: upload kaggle.json, permissions, download + unzip dataset.
    \item Verified counts per split and label folder; validated directory structure.
    \item Built file path lists using os.walk; shuffled; confirmed lengths.
    \item Preprocessing and augmentation with ImageDataGenerator:
      rescale, rotation, shifts, shear, zoom, horizontal flip, fill mode.
    \item Generators via flow\_from\_directory with target\_size=(224,224), batch\_size=32, class\_mode=binary.
    \item Visual sanity checks: plotting sample batches from train/val/test generators.
  \end{itemize}
\item Modeling work (baseline + transfer learning + fine-tuning):
  \begin{itemize}[leftmargin=*,nosep]
    \item Baseline custom CNN (Functional API): Conv2D blocks (16/32/64), BatchNorm, ReLU, MaxPool, Dropout, Flatten, Dense, Sigmoid.
    \item Callbacks: ModelCheckpoint (best weights), EarlyStopping, ReduceLROnPlateau.
    \item Transfer learning: ResNet152V2 (ImageNet, include\_top=False), frozen base, custom head (GAP, Dense, BatchNorm, Dropout, Sigmoid).
    \item Fine-tuning: unfreeze base model with controlled layer freezing (e.g., freeze all but last ~30 layers), careful LR settings, monitoring overfitting.
  \end{itemize}
\item Deployment workflow (Hugging Face + Gradio):
  \begin{itemize}[leftmargin=*,nosep]
    \item Hugging Face login (CLI/notebook), repository use, versioning with add/commit/push concept.
    \item Model saving (.h5/.keras) and packaging for inference.
    \item Gradio demo: image upload $\rightarrow$ resize to 224x224 $\rightarrow$ normalize $\rightarrow$ predict $\rightarrow$ label (Pneumonia/Normal) with threshold 0.5.
    \item Validated end-to-end inference usability.
  \end{itemize}
\item Final presentation preparation: problem statement, dataset context, preprocessing, models, evaluation approach, deployment screenshots/steps, and lessons learned.
}
{
Completed an end-to-end ML project under an intensive schedule: dataset ingestion, preprocessing, baseline CNN, transfer learning + fine-tuning, and deployment using Hugging Face/Gradio, supported by MLOps and prompt engineering practices.
}
{
\item Chest X-ray pneumonia dataset structure and clinical context
\item Keras data generators and augmentation for images
\item Custom CNN vs transfer learning vs fine-tuning
\item Callbacks for training control and reproducibility
\item Hugging Face Hub/Spaces deployment and Gradio UI
\item Prompt engineering patterns for debugging, documentation, evaluation
}
{
End-to-end ML project delivery from data to deployment
}
{
\item Problem-solving Skills
\item Quantitative \& Analytical Skills
}
{
\item Leadership Skills
\item Communication Skills
}

%========================================
% Helpful Links (Credible)
%========================================
\newpage
{\Large \textbf{Helpful References (Links)}}\\[6pt]
\begin{itemize}[leftmargin=*]
  \item Kaggle dataset (Chest X-Ray Pneumonia): \url{https://www.kaggle.com/datasets/paultimothymooney/chest-xray-pneumonia}
  \item Dataset clinical context / paper page (Cell): \url{http://www.cell.com/cell/fulltext/S0092-8674(18)30154-5}
  \item TensorFlow ImageDataGenerator docs: \url{https://www.tensorflow.org/api_docs/python/tf/keras/preprocessing/image/ImageDataGenerator}
  \item TensorFlow Transfer Learning guide: \url{https://www.tensorflow.org/tutorials/images/transfer_learning}
  \item Keras Callbacks (EarlyStopping / ReduceLROnPlateau / ModelCheckpoint): \url{https://keras.io/api/callbacks/}
  \item Hugging Face main site: \url{https://huggingface.co/}
  \item Hugging Face Hub docs: \url{https://huggingface.co/docs/hub/index}
  \item Hugging Face Transformers docs: \url{https://huggingface.co/docs/transformers/index}
  \item Hugging Face Spaces docs: \url{https://huggingface.co/docs/hub/spaces}
  \item Gradio documentation: \url{https://www.gradio.app/docs}
  \item MLflow official documentation: \url{https://mlflow.org/docs/latest/index.html}
\end{itemize}

\end{document}
